%% !TEX program = xelatex
\documentclass[12pt]{article}

% Поддержка кириллицы и шрифтов
\usepackage{fontspec}
\usepackage[russian]{babel}
\usepackage{csquotes}

\usepackage{tabularx}
\usepackage[table]{xcolor} % ← для цвета в таблицах

% Библиография (стиль ГОСТ)
\usepackage[backend=biber, style=gost-numeric]{biblatex}
\addbibresource{references.bib} % создайте этот файл отдельно

% Шрифт по ГОСТ
\setmainfont{Times New Roman}

% Поля страницы (по умолчанию ~3 см; для точного ГОСТ используйте пакет geometry)
\usepackage[a4paper, margin=2.5cm]{geometry}

% Заголовки разделов без номеров (как в большинстве журналов)
\usepackage{titlesec}
\titleformat{\section}{\normalfont\bfseries\large}{\thesection}{1em}{}
\titlespacing*{\section}{0pt}{12pt}{6pt}

\title{Название научной статьи}
\author{Самаркин А.И.\\
        \small Научная организация, Город}
\date{}

\begin{document}

\maketitle

% Аннотация
\begin{abstract}
\noindent\textbf{Аннотация.} Краткое изложение цели, методов, основных результатов и выводов исследования (150–250 слов). Аннотация должна быть информативной, содержать ключевые слова и отражать суть работы без дополнительных пояснений. Используется сплошной текст без переносов и сокращений.
\end{abstract}

% Ключевые слова
\noindent\textbf{Ключевые слова:} первое ключевое слово, второе, третье, четвёртое, пятое.

% Основной текст
\section*{Введение}
Обоснование актуальности темы исследования. Краткий обзор литературы с выявлением <<пробела>> в знаниях, который закрывает данная работа. Чёткая формулировка цели и задач исследования. Гипотеза (если применимо). Объём — 15–20\% от общего текста.

Пример ссылки на источник~\cite{ivanov2020}. Теоретическая или практическая значимость работы.

\section*{Методы}
Описание объекта и методов исследования с достаточной детализацией для воспроизведения эксперимента другими исследователями. Указание оборудования, программного обеспечения, статистических методов, этических аспектов (если применимо).

Пример: анализ данных проводился с использованием пакета \textit{SciPy} (версия 1.12)~\cite{scipy2020}.

\section*{Результаты}
Представление полученных данных без интерпретации. Использование таблиц, рисунков, графиков с чёткой нумерацией и подписями. Статистические показатели (среднее, доверительные интервалы, p-уровни).

Пример: как показано на рис.~1, наблюдается достоверное различие между группами ($p < 0{,}01$).

\section*{Обсуждение}
Интерпретация результатов: сравнение с данными других авторов, объяснение неожиданных находок, обсуждение ограничений исследования. Ответ на вопрос, поставленный во Введении. Формулировка выводов и предложений для будущих исследований.

Наши данные согласуются с результатами Петрова~\cite{petrov2019}, однако противоречат выводам Смирнова~\cite{smirnov2021}.

\section*{Заключение}
Краткие выводы (3–5 пунктов), отражающие достижение цели исследования. Практические рекомендации (если есть).

\printbibliography[title={Список литературы}]

% Примеры рисунков и таблиц (раскомментируйте при необходимости)
\begin{figure}[ht]
\centering
\fbox{Здесь будет график}
\caption{Подпись к рисунку}
\end{figure}

%
\begin{table}[ht]
\centering
\caption{Пример таблицы с вертикальными линиями}
\begin{tabular}{|l|c|c|}
\hline
\textbf{Показатель} & \textbf{Группа А} & \textbf{Группа Б} \\
\hline
Среднее значение    & 12.5             & 18.3             \\
Стандартное отклонение & 2.1           & 3.4              \\
\hline
\end{tabular}
\end{table}

\begin{table}[ht]
\caption{Пример таблицы с цветным заголовком}
\begin{tabular}{|l|c|c|}
\hline
\rowcolor{yellow}
\textbf{Параметр} & \textbf{Значение 1} & \textbf{Значение 2} \\
\hline
Температура       & 25 °C              & 37 °C              \\
Время реакции     & 10 мин             & 15 мин             \\
\hline
\end{tabular}
\end{table}

\end{document}